% Independent internal foundation fragment, 2026-11-10.
% Requires amsmath, amssymb, amsthm, and theorem/lemma environments.
% The surrounding development supplies P,Z,G,L,J, the numerical equality
% test delta, their endpoint laws, and derived quantifier-free induction.
% No primitive unrestricted addition, order, or constructor-coverage axiom
% is assumed in this fragment.

\section{Internal order, difference, and addition under a bound}
\label{sec:internal-order-foundations}

This attributed appendix supplies the arithmetic facts needed by the bounded compiler~\cite{bounded}. The notation $\widehat d$ distinguishes its constructed difference from the primitive $d$ of our new basis.
All its inductions are on quantifier-free formulas with the other displayed
variables held fixed as parameters.  In particular, an induction below is
never shorthand for induction on a universally quantified formula.  Use $\epsilon$ for the numerical equality test, with values zero and one, and use
the already derived terms and equations
\[
 P0=0,\quad P(Sx)=x,\qquad Z0=1,\quad Z(Sx)=0,
 \qquad G(0,a)=0,\quad G(Sx,a)=a,
\]
\[
 L_h(0)=0,\quad L_h(Sn)=h(n,L_h(n)),\qquad
 J(a,0)=a,\quad J(a,1)=Sa.
\]
Here $1=S0$.  Applying $P$ to an equality of successors gives successor
injectivity.  Applying $Z$ to $0=Sx$ contradicts $0\ne1$, so zero is
distinct from every successor.  The proved Boolean range and reflection laws consequently give
\[
 \epsilon(x,x)=1,\quad \epsilon(x,y)=\epsilon(y,x),\quad
 \epsilon(Sx,Sy)=\epsilon(x,y),\quad \epsilon(0,Sx)=0.
\]
Each identity follows by the two possible Boolean values and the reflection
law for $\epsilon$; these are derived identities, not extra axioms.

\begin{lemma}[Derived constructor coverage]
\label{lem:internal-coverage}
The theory proves $x=0\lor x=S(Px)$.
\end{lemma}
\begin{proof}
Apply the already derived quantifier-free induction to
$x=0\lor x=S(Px)$.  The base is immediate.  At a successor the second
disjunct follows directly from $P(Sx)=x$, independently of the induction
hypothesis.  Thus constructor coverage need not be a primitive axiom.
\end{proof}

\subsection{A hit term defines a discrete linear order}

Define
\[
 H(n,y):=L(n;r,v.J(v,\epsilon(r,y))).
\]
Its defining equations are
\[
 H(0,y)=0,\qquad H(Sn,y)=J(H(n,y),\epsilon(n,y)).
 \tag{H}\label{eq:internal-hit-recursion}
\]

\begin{lemma}[Hit identities and duality]
\label{lem:internal-hit}
The following identities and range assertion hold internally:
\begin{align*}
 H(Sn,0)&=1,& H(Sn,Sy)&=H(n,y),& H(n,n)&=0,\\
 H(Sn,n)&=1,& H(n,y)&=0\lor H(n,y)=1,\\
 H(Sy,x)&=Z(H(x,y)).
\end{align*}
Moreover a hit persists: $H(n,y)=1$ implies $H(Sn,y)=1$.
\end{lemma}
\begin{proof}
Equation induction on $n$ proves $H(Sn,0)=1$: its first step hits zero,
and every later equality test against zero is zero.  Equation induction
on $n$, with $y$ fixed, proves $H(Sn,Sy)=H(n,y)$, using
$\epsilon(Sn,Sy)=\epsilon(n,y)$.  A further equation induction gives
$H(n,n)=0$; the recurrence then gives $H(Sn,n)=1$.

For Boolean range use quantifier-free induction on $n$, with $y$ fixed.
If $\epsilon(n,y)=0$, the state is unchanged.  If $\epsilon(n,y)=1$, reflection
gives $n=y$, so the old state is $H(n,n)=0$ and the new state is one.
If $H(n,y)=1$, diagonal zero shows $n\ne y$, so the next equality test
vanishes.  This proves persistence.  The same case split shows
\[
 H(Sn,y)=1\ \Longrightarrow\ H(n,y)=1\ \lor\ n=y.
 \tag{New hit}\label{eq:internal-new-hit}
\]

Finally prove duality by equation induction on $x$, with $y$ fixed.
At zero both sides are one.  In the successor step the new left side is
$H(Sy,Sx)=H(y,x)$, whereas the old left side is
$J(H(y,x),\epsilon(y,x))$.  If $\epsilon(x,y)=0$, both relevant recurrences
leave their states unchanged, and the induction hypothesis gives the
required equality.  If $\epsilon(x,y)=1$, put $x=y$; the new left side is
zero and the new right side is $Z(H(Sx,x))=Z1=0$.
\end{proof}

Put
\[
 x<y\quad:\!\Longleftrightarrow\quad H(y,x)=1,
 \qquad
 x\le y\quad:\!\Longleftrightarrow\quad H(Sy,x)=1.
 \tag{Order}\label{eq:internal-order}
\]
These are quantifier-free abbreviations.

\begin{lemma}[Discrete linear-order laws]
\label{lem:internal-linear-order}
The relations in \eqref{eq:internal-order} form a discrete linear order
with least element zero.  In particular,
\[
 x\le y\ \Longleftrightarrow\ (x=y\lor x<y),\qquad
 x<Sy\ \Longleftrightarrow\ x\le y,
\]
successor preserves and reflects both comparisons, and
$x\le Sy$ with $x\ne Sy$ implies $x\le y$.
\end{lemma}
\begin{proof}
Diagonal identities give reflexivity of $\le$, irreflexivity of $<$,
and $x<Sx$.  The recurrence for $H(Sy,x)$, split on $\epsilon(y,x)$,
gives the first displayed equivalence.  Duality and Boolean range give
trichotomy: if $x\ne y$, then $H(Sy,x)=H(y,x)$, and $H(x,y)$ is its
Boolean complement.  Thus exactly one of $x<y$, $x=y$, $y<x$ holds.

For strict transitivity retain $H(y,x)=1$ and induct on $z$ in the
quantifier-free formula
\[
                 H(z,y)=1\ \longrightarrow\ H(z,x)=1.
\]
The base has false antecedent.  At the successor, the new-hit implication
\eqref{eq:internal-new-hit} gives either $H(z,y)=1$, when the induction
hypothesis followed by persistence applies, or $z=y$, when the retained
$H(y,x)=1$ followed by persistence applies.  Transitivity of $\le$ and
the mixed comparison laws now follow by splitting equality from strict
comparison.  Trichotomy gives antisymmetry.

The identities $H(Sn,Sy)=H(n,y)$ give successor preservation and
reflection.  The equivalence $x<Sy\leftrightarrow x\le y$ is immediate
from the definitions, and the asserted discrete predecessor step follows
from it and the first equivalence.  Finally $H(Sx,0)=1$ gives $0\le x$;
if $x\le0$, trichotomy or constructor coverage gives $x=0$.
\end{proof}

\subsection{Difference is an internal term, not an imported operation}

Define
\[
 \widehat d(x,y):=L(x;u,v.J(v,H(Su,y))).
\]
Thus
\[
 \widehat d(0,y)=0,
 \qquad \widehat d(Sx,y)=J(\widehat d(x,y),H(Sx,y)).
 \tag{d}\label{eq:internal-difference-recursion}
\]

\begin{lemma}[Difference laws]
\label{lem:internal-difference}
The theory proves
\begin{align*}
 \widehat d(x,0)&=x,& \widehat d(Sx,Sy)&=\widehat d(x,y),&\widehat d(x,x)&=0,\\
 \widehat d(x,y)=0&\ \Longleftrightarrow\ x\le y,
 &\widehat d(x,Sy)&=P(\widehat d(x,y)),&\widehat d(x,y)&\le x.
\end{align*}
Difference is monotone in its first argument and antitone in its second.
Furthermore
\begin{align*}
 y\le x&\ \Longrightarrow\ \widehat d(Sx,y)=S(\widehat d(x,y)),\\
 0<y\le x&\ \Longrightarrow\ \widehat d(x,Py)=S(\widehat d(x,y)).
\end{align*}
\end{lemma}
\begin{proof}
Equation induction gives $\widehat d(x,0)=x$ because $H(Sx,0)=1$.
For simultaneous successors, induct on $x$ with $y$ fixed.  The base
uses $H(S0,Sy)=0$.  At the successor use
$H(S(Sx),Sy)=H(Sx,y)$ and the induction hypothesis in
\eqref{eq:internal-difference-recursion}.

Quantifier-free induction on $x$, with $y$ fixed, proves
$\widehat d(x,y)=0\leftrightarrow H(x,y)=0$.
At the successor, if $H(Sx,y)=1$, the new difference is a successor and
is nonzero.  If $H(Sx,y)=0$, persistence implies $H(x,y)=0$, so the old
difference is zero and remains zero.  Duality identifies $H(x,y)=0$
with $H(Sy,x)=1$, proving the order characterization and $\widehat d(x,x)=0$.

Now induct on $x$ for $\widehat d(x,Sy)=P(\widehat d(x,y))$, with $y$ fixed.  Its successor
left side is $\widehat d(x,y)$.  The right side is
$P(J(\widehat d(x,y),H(Sx,y)))$.  When the Boolean is one, predecessor cancels
the successor; when it is zero, the preceding zero characterization
gives $\widehat d(x,y)=0$.  Both cases give the desired equality.

Quantifier-free induction on $x$ proves $\widehat d(x,y)\le x$: its successor
state is either the old state or its successor.  For first-argument
monotonicity induct on $v$ in
\[
       u\le v\ \longrightarrow\ \widehat d(u,z)\le \widehat d(v,z),
\]
with $u,z$ fixed.  At $Sv$, either $u=Sv$ or discreteness gives $u\le v$;
in the latter case use the induction hypothesis and the fact that the
new state is at least its old state.  The base uses $u\le0\to u=0$.

For second-argument antitonicity retain $y\le z$ and induct on $x$ in
$\widehat d(x,z)\le \widehat d(x,y)$.  If $H(Sx,z)=0$, the left state stays fixed and
the right state cannot decrease.  If $H(Sx,z)=1$, then $z\le x$;
transitivity gives $y\le x$, so both states increment and successor
preserves their old comparison.

The first final implication is the defining recurrence, since
$y\le x$ means $H(Sx,y)=1$.  For the second, coverage gives
$y=S(Py)$ when $y>0$.  The inequalities imply $Py<x$, so
$\widehat d(x,Py)\ne0$.  Functional coverage at this positive difference and
$\widehat d(x,S(Py))=P(\widehat d(x,Py))$ yield
$\widehat d(x,Py)=S(\widehat d(x,y))$.
\end{proof}

Define
\[
                      \operatorname{mn}(x,y):=\widehat d(x,\widehat d(x,y)).
\]

\begin{lemma}[Minimum]
\label{lem:internal-minimum}
If $x\le y$, then $\operatorname{mn}(x,y)=x$; if $y\le x$, then $\operatorname{mn}(x,y)=y$.
Consequently $\operatorname{mn}$ is the minimum operation for the internal linear order.
\end{lemma}
\begin{proof}
The first claim uses $\widehat d(x,y)=0$ and $\widehat d(x,0)=x$.
For the second induct on $y$ in
$y\le x\to \widehat d(x,\widehat d(x,y))=y$, with $x$ fixed.
The base is $\widehat d(x,x)=0$.  If $Sy\le x$, then $y<x$ and
$0<\widehat d(x,y)\le x$.  Apply the final difference implication to obtain
\[
 \widehat d(x,\widehat d(x,Sy))
   =\widehat d(x,P(\widehat d(x,y)))
   =S(\widehat d(x,\widehat d(x,y)))=Sy.
\]
The last equality uses the induction hypothesis, which applies because
$y\le x$.  Linearity completes the interpretation of $\operatorname{mn}$.
\end{proof}

\subsection{Addition with an explicit common bound}

Define the unary doubling term
\[
             \operatorname{db}(n):=L(n;u,v.S(Sv)).
\]
Thus $\operatorname{db}(0)=0$ and $\operatorname{db}(Sn)=S(S(\operatorname{db}(n)))$.  The notation $2n$ in the compiler
may be read as this term $\operatorname{db}(n)$; it does not invoke binary addition.

\begin{lemma}[Doubling identities]
\label{lem:internal-doubling}
The term $\operatorname{db}$ is monotone, $n\le \operatorname{db}(n)$, and $\widehat d(\operatorname{db}(n),n)=n$.
\end{lemma}
\begin{proof}
Quantifier-free induction gives $n\le \operatorname{db}(n)$ directly from the two
successors in the defining equation.  Monotonicity follows by induction
on the larger argument, splitting equality at its successor from the
discrete predecessor case, just as for first-argument monotonicity of $d$.
Equation induction proves the remaining identity.  At the successor,
\[
 \widehat d(\operatorname{db}(Sn),Sn)
   =\widehat d(S(S(\operatorname{db}(n))),Sn)=\widehat d(S(\operatorname{db}(n)),n)
   =S(\widehat d(\operatorname{db}(n),n))=Sn,
\]
where the third equality uses $n\le \operatorname{db}(n)$.
\end{proof}

Put
\[
 K_b(a):=\widehat d(\operatorname{db}(b),\widehat d(b,a)),\qquad
 \operatorname{ad}_b(a,c):=\widehat d(K_b(a),\widehat d(b,c)).
 \tag{Bounded addition}\label{eq:internal-bounded-addition}
\]
Every expression here is a finite term of the original language.

\begin{lemma}[Bounded-addition laws]
\label{lem:internal-bounded-addition}
Internally, the following hold:
\begin{enumerate}
\item If $a\le b$, then $\operatorname{ad}_b(a,0)=a$.
\item If $a\le b$ and $Sc\le b$, then
      $\operatorname{ad}_b(a,Sc)=S(\operatorname{ad}_b(a,c))$.
\item If $c\le b$, then $\operatorname{ad}_b(0,c)=c$.
\item $\operatorname{ad}_b(b,b)=\operatorname{db}(b)$ and $\operatorname{ad}_b(a,c)\le \operatorname{db}(b)$ for all $a,c$.
\item $\operatorname{ad}_b$ is monotone in its second argument.
\item If $a,c\le b$ and $a,c\le b'$, then
      $\operatorname{ad}_b(a,c)=\operatorname{ad}_{b'}(a,c)$.
\end{enumerate}
Thus it supplies addition on inputs bounded by $b$, without asserting an
unrestricted binary addition term.
\end{lemma}
\begin{proof}
Since $\widehat d(b,a)\le b$, antitonicity and doubling give
\[
            b=\widehat d(\operatorname{db}(b),b)\le K_b(a)\le \operatorname{db}(b).
 \tag{K bounds}\label{eq:internal-K-bounds}
\]
Also $K_b(0)=b$.  If $Sa\le b$, put $q=\widehat d(b,a)$.
Then $0<q\le b\le \operatorname{db}(b)$ and $\widehat d(b,Sa)=Pq$, so
\[
                  K_b(Sa)=\widehat d(\operatorname{db}(b),Pq)=S(K_b(a)).
\]
Quantifier-free induction on $a$ now proves the first assertion under
$a\le b$.  The base is $\widehat d(b,b)=0$; at the successor use the last
identity and $b\le K_b(a)$ to increment $\widehat d(K_b(a),b)$.

For the second assertion put $q=\widehat d(b,c)$.  The premise $Sc\le b$ gives
$0<q\le b\le K_b(a)$.  Hence
\[
 \operatorname{ad}_b(a,Sc)
   =\widehat d(K_b(a),Pq)=S(\widehat d(K_b(a),q))=S(\operatorname{ad}_b(a,c)).
\]
The third assertion follows by quantifier-free induction on $c\le b$
using the first two.  For the fourth, substitute $a=c=b$ and use
$\widehat d(b,b)=0$ twice; the upper bound follows from
$\widehat d(K_b(a),\widehat d(b,c))\le K_b(a)\le \operatorname{db}(b)$.

If $c\le c'$, antitonicity gives $\widehat d(b,c')\le \widehat d(b,c)$; applying
antitonicity again to the outer difference gives
$\operatorname{ad}_b(a,c)\le \operatorname{ad}_b(a,c')$.  Finally retain $a\le b,b'$ and induct on $c$
in the quantifier-free implication
\[
 c\le \operatorname{mn}(b,b')\ \longrightarrow\ \operatorname{ad}_b(a,c)=\operatorname{ad}_{b'}(a,c).
\]
The base gives $a$ on each side, and whenever $Sc\le \operatorname{mn}(b,b')$ the
successor laws apply to both sides and preserve the induction hypothesis.
\end{proof}

These proofs establish the numerical prerequisites for bounded selection
and clipped recursion entirely within the proposed theory.  Each use of
a case split is justified by numerical equality decidability, Boolean
range, or the derived linear-order and coverage laws.  No standard-model
correctness assertion and no hidden unrestricted addition operation is
used as an internal axiom.

\begin{corollary}[Identification with the new primitive monus]
In $T_{Bd}$, $\widehat d(x,y)=d(x,y)$ for every $x,y$.
\end{corollary}
\begin{proof}
Keep $x$ fixed and apply literal-equation induction in $y$.
Both terms have value $x$ at zero. Their successor values are the
predecessors of their previous values, by the proved difference lemma
and the primitive monus equation. Congruence of $P$ proves the step.
All occurrences of $\epsilon$ in $\widehat d$ are inlined as the fixed
$B+d$ term from the main text. Thus this is a source-language equation
induction, without an appeal to a graph-function extension.
\end{proof}
